% =====================================================================
% Background and Literature Review  --  TARGET: ~600 words
% Marking weight: 10 marks (combined with abstract + intro)
% =====================================================================
\section{Background and Literature Review}
\label{sec:literature}

\guide{
The marking rubric for 1st+ wants \textbf{critical analysis with deep
insights and strong academic references}. Aim for $\geq$10 cited
sources spread across the four sub-topics below. Do not just list
papers; \emph{compare} them and \emph{say what your project takes from
each}.

The four sub-topics map directly onto the brief:
\begin{enumerate}
  \item Industrial automation \& multi-DOF arms (need / DOF theory /
        ISO~9283 metrics).
  \item Wireless protocols (BT Classic vs BLE vs Zigbee vs Wi-Fi vs
        nRF24): include a trade-off table.
  \item Microcontroller selection (Arduino vs Pi vs ESP32) — this is
        where you justify rejecting the Pi for real-time control.
  \item Closed-loop control \& feedback sensing (PID, potentiometer
        vs encoder vs force sensor).
\end{enumerate}
References~\cite{ifr2024,siciliano2009,spong2020,astrom2008,iso9283,
ieee802151,ieee802154,gomez2012,baker2005,atmega2560} are already
seeded in references.bib --- use them and add 3--4 more recent
(2022--2025) papers per sub-topic for currency.
}

\subsection{Articulated manipulators and degrees of freedom}
\placeholder{[\ldots one paragraph on industrial-arm DOF theory,
citing Siciliano~\cite{siciliano2009} and Spong~\cite{spong2020};
explain why 4-DOF + gripper is the natural minimum for sort tasks
and matches the brief\ldots]}

\subsection{Performance metrics for educational manipulators}
\placeholder{[\ldots one paragraph introducing ISO~9283 repeatability
and accuracy metrics~\cite{iso9283}; contrast industrial sub-mm
performance with what's achievable with hobby servos; explain that
closing the loop on potentiometer feedback is the standard
mitigation\ldots]}

\subsection{Wireless protocols for embedded teleoperation}
\guide{This is the strongest place to demonstrate critical analysis.
Compare at least four protocols. The structure that scores well:
\begin{itemize}
  \item One sub-paragraph per protocol (Bluetooth Classic, BLE,
        Zigbee, Wi-Fi). For each, note: data rate, latency, power,
        topology limit, and one named drawback.
  \item Closing paragraph: explain why Bluetooth Classic via HC-05
        was chosen, citing Baker's comparison~\cite{baker2005} and
        Gomez et al.~\cite{gomez2012}.
  \item Mention 2.4\,GHz interference and adaptive frequency hopping.
\end{itemize}
A trade-off table here picks up easy marks.}

\placeholder{[\ldots three to four short sub-paragraphs on the four
protocols, then a justification paragraph for HC-05\ldots]}

\subsection{Microcontroller selection for real-time control}
\guide{This is where the Pi-vs-Arduino debate goes. Key bullets:
non-deterministic Linux scheduling, no native ADC, long boot, high
idle current. Cite the ATmega datasheet~\cite{atmega2560}. Land on
the compromise: Arduino for real-time motor control, Pi (or ESP32)
as a vision \emph{co-processor} only.}

\placeholder{[\ldots one paragraph rejecting the Pi for hard real-time
control, then justifying Arduino Mega for motor/feedback work and
introducing the heterogeneous architecture used here\ldots]}

\subsection{Closed-loop joint control}
\guide{Short. Cite \AA str\"om \& Murray~\cite{astrom2008}; mention
discrete-time PID at 50--200\,Hz; contrast potentiometer feedback
(absolute, cheap, low-resolution) with encoders (relative, expensive,
high-resolution) and force sensors. Set up
Section~\ref{sec:closedloop} to do the implementation detail.}

\placeholder{[\ldots one short paragraph on PID for servo joints,
ending by forward-referencing the implementation section\ldots]}
